\hypertarget{advanced-data-structures-internals}{%
\chapter{Advanced Data Structures
Internals}\label{advanced-data-structures-internals}}

While Python's \passthrough{\lstinline!list!} and
\passthrough{\lstinline!dict!} are versatile (covered in Chapter 2 and
5), the \passthrough{\lstinline!collections!} and related modules
provide specialized structures optimized for specific algorithmic
complexities. Understanding their C implementations is key to writing
high-performance code.

\hypertarget{collections.deque-the-doubly-linked-block-architecture}{%
\subsection{\texorpdfstring{33.1 \texttt{collections.deque}: The
Doubly-Linked Block
Architecture}{33.1 collections.deque: The Doubly-Linked Block Architecture}}\label{collections.deque-the-doubly-linked-block-architecture}}

A \passthrough{\lstinline!list!} is an array-based structure, making
insertions/deletions at the start \(O(N)\). A
\passthrough{\lstinline!deque!} (Double-Ended Queue) is designed for
\(O(1)\) appends and pops from both ends.

\hypertarget{the-block-structure}{%
\subsubsection{1. The Block Structure}\label{the-block-structure}}

Unlike a standard doubly-linked list where each node holds one element
(high memory overhead), CPython's \passthrough{\lstinline!deque!} uses a
\textbf{doubly-linked list of blocks}. * Each block is a fixed-size
array (typically 64 elements). * The \passthrough{\lstinline!deque!}
object tracks the \passthrough{\lstinline!leftblock!},
\passthrough{\lstinline!rightblock!}, and indices for the start and end.

\hypertarget{performance-implications}{%
\subsubsection{2. Performance
Implications}\label{performance-implications}}

\begin{itemize}
\tightlist
\item
  \textbf{Memory Efficiency}: By grouping elements into blocks, it
  minimizes the number of \passthrough{\lstinline!malloc!} calls and
  improves cache locality compared to a simple linked list.
\item
  \textbf{No Reallocations}: Unlike \passthrough{\lstinline!list!},
  which may need to \passthrough{\lstinline!realloc!} and copy the
  entire array, a \passthrough{\lstinline!deque!} simply allocates a new
  block when the current one is full, making its performance extremely
  predictable for streaming data.
\end{itemize}

\hypertarget{heapq-binary-heaps-on-arrays}{%
\subsection{\texorpdfstring{33.2 \texttt{heapq}: Binary Heaps on
Arrays}{33.2 heapq: Binary Heaps on Arrays}}\label{heapq-binary-heaps-on-arrays}}

The \passthrough{\lstinline!heapq!} module provides a min-heap
implementation using a standard Python \passthrough{\lstinline!list!}.

\hypertarget{min-heap-invariant}{%
\subsubsection{1. Min-Heap Invariant}\label{min-heap-invariant}}

For every node at index \passthrough{\lstinline!i!}, its children are at
\passthrough{\lstinline!2*i + 1!} and \passthrough{\lstinline!2*i + 2!}.
The value at \passthrough{\lstinline!i!} is always less than or equal to
its children.

\hypertarget{the-_siftdown-and-_siftup-mechanics}{%
\subsubsection{\texorpdfstring{2. The \texttt{\_siftdown} and
\texttt{\_siftup}
Mechanics}{2. The \_siftdown and \_siftup Mechanics}}\label{the-_siftdown-and-_siftup-mechanics}}

When you \passthrough{\lstinline!heappush!}, the element is appended to
the list and then ``sifted up'' (swapped with parents) until the
invariant is restored. \passthrough{\lstinline!heappop!} replaces the
root with the last element and ``sifts down.'' Both operations are
\(O(\log N)\).

\hypertarget{bisect-binary-search-optimization}{%
\subsection{\texorpdfstring{33.3 \texttt{bisect}: Binary Search
Optimization}{33.3 bisect: Binary Search Optimization}}\label{bisect-binary-search-optimization}}

The \passthrough{\lstinline!bisect!} module implements binary search on
sorted sequences. * \textbf{\passthrough{\lstinline!bisect\_left!} /
\passthrough{\lstinline!bisect\_right!}}: Find the insertion point for
an element to maintain order. * \textbf{Internals}: These are
implemented in C for speed. They perform a simple \(O(\log N)\)
bisection, assuming the underlying sequence supports random access
(\(O(1)\) indexing).

\hypertarget{collections.counter-and-defaultdict}{%
\subsection{\texorpdfstring{33.4 \texttt{collections.Counter} and
\texttt{defaultdict}}{33.4 collections.Counter and defaultdict}}\label{collections.counter-and-defaultdict}}

These are thin wrappers around the standard
\passthrough{\lstinline!dict!}. *
\textbf{\passthrough{\lstinline!defaultdict!}}: Overrides
\passthrough{\lstinline!\_\_missing\_\_!} (a dunder method called by
\passthrough{\lstinline!dict.\_\_getitem\_\_!} when a key isn't found).
It calls the \passthrough{\lstinline!default\_factory!} and inserts the
result. * \textbf{\passthrough{\lstinline!Counter!}}: Primarily adds the
\passthrough{\lstinline!most\_common()!} method (which uses a
\passthrough{\lstinline!heapq!} for large \(K\) or sorting for small
\(K\)) and operator overloading for set-like addition/subtraction.

\hypertarget{weakref-memory-management-proxies}{%
\subsection{\texorpdfstring{33.5 \texttt{weakref}: Memory Management
Proxies}{33.5 weakref: Memory Management Proxies}}\label{weakref-memory-management-proxies}}

Normally, assigning an object to a variable increases its reference
count. A \passthrough{\lstinline!weakref!} allows you to reference an
object \textbf{without} increasing its reference count.

\hypertarget{the-weakref.ref-object}{%
\subsubsection{\texorpdfstring{1. The \texttt{weakref.ref}
Object}{1. The weakref.ref Object}}\label{the-weakref.ref-object}}

A weak reference is a small proxy object. When the referent's reference
count drops to zero, the referent is garbage collected, and the weakref
is automatically set to \passthrough{\lstinline!None!}.

\hypertarget{internal-callbacks}{%
\subsubsection{2. Internal Callbacks}\label{internal-callbacks}}

You can register a callback that executes immediately when the referent
is about to be destroyed. This is used extensively in internal caches
(like \passthrough{\lstinline!functools.lru\_cache!} for objects) to
prevent memory leaks in long-running processes.

\begin{center}\rule{0.5\linewidth}{0.5pt}\end{center}

\begin{center}\rule{0.5\linewidth}{0.5pt}\end{center}
